\documentclass{article}
\usepackage{amsmath,amssymb,amsthm}
\usepackage{microtype}
\usepackage{hyperref}
\usepackage[margin=1in]{geometry}

\newtheorem{theorem}{Theorem}
\newtheorem{proposition}[theorem]{Proposition}
\newtheorem{lemma}[theorem]{Lemma}
\newtheorem{corollary}[theorem]{Corollary}
\newtheorem{conjecture}[theorem]{Conjecture}
\theoremstyle{remark}
\newtheorem{remark}[theorem]{Remark}

\newcommand{\R}{\mathbb{R}}
\newcommand{\E}{\mathbb{E}}
\newcommand{\tr}{\mathrm{tr}}
\newcommand{\rank}{\mathrm{rank}}
\newcommand{\range}{\mathrm{range}}
\newcommand{\cheb}{R_{\mathrm{c}}}
\newcommand{\deff}{d_{\mathrm{eff}}}
\newcommand{\tauH}{\bar\tau_{H}}
\newcommand{\Hbar}{\bar H}
\newcommand{\Stiefel}{\mathrm{St}}
\newcommand{\hpinv}{^{+}}

\title{Rank-$r$ Rate-Distortion Function for Model Merging
  (v1 scaffold, Phase 2 Day 14)}
\author{Sankalp Pathak}
\date{April 2026}

\begin{document}
\maketitle

\begin{abstract}
Phase 1 scaffold. We lift the Phase 0 toy theorem
(\texttt{theory/toy\_theorem\_v0.tex}, $H_t = I_d$ isotropic) to the
rank-$r$ weight-disentangled setting of Ortiz-Jimenez et~al.\
(NeurIPS~2023), where each task loss is a quadratic with Hessian
$H_t$ supported on an $r$-dimensional subspace $V_t \subset \R^d$.
The Phase 0 cross-term-vanishing identity generalizes to an
$H_t$-weighted centroid $\tauH$ (\S\ref{sec:id}). Under a
Stiefel-random hard distribution $P$ (\S\ref{sec:hard}), the floor
of that identity has a closed form in $\deff := \rank(\sum_t P_{V_t})$
(Lemma~\ref{lem:floor-rankr}, \S\ref{sec:floor}), and for
$H_t = P_{V_t}$ the rank-$r$ rate--distortion lower bound
\[
  \mathcal{D}^\star \;\geq\; B^2\bigl(1-\E_P[\deff]/(Tr)\bigr)
  + c_{\mathrm{TQ}}\,(B^2\E_P[\deff]/(Tr))\,2^{-2R/\deff}
\]
is proved (Theorem~\ref{thm:rankr}, \S\ref{sec:thm}). Three
sanity-check limits (\S\ref{sec:limits}) fall out as corollaries by
substitution of $\deff$: (i) $r=d$ recovers Phase 0 verbatim; (ii)
$T=1$ recovers TurboQuant~Theorem~3 on a rank-$r$ source; (iii)
shared subspace $V_t = V$ recovers Phase 0 with $d \mapsto r$.
\textbf{Day-10/11 extension (\S\ref{sec:general-ht}):} the same bound
holds for general $H_t \succeq 0$ with $\range(H_t) \subseteq V_t$
and \emph{arbitrary task-dependent} spectra $(D_t)_{t=1}^T$
(Theorem~\ref{thm:general}), under the generalized hard
distribution $P^\star$ where $\tau_t$ is uniform on the $H_t$-
ellipsoid in $V_t$. This closes the $\text{MSE}\to\text{CE}$
bridge for LoRA Fisher metrics, including heterogeneous curvature
across tasks.
A matching \textbf{achievability} upper bound
(Theorem~\ref{thm:achv}, \S\ref{sec:achv}) via an explicit
Gaussian-QR random-orthogonal mixer plus uniform scalar
quantization of $\eta := \bar H^{1/2}\tauH$ closes the floor-zero
Stiefel regime up to a universal constant $C = O(T)$, pinning
down $\mathcal{D}^\star = \Theta(B^2\cdot 2^{-2R/(Tr)})$.
Numerical support in \texttt{code/synthetic/day8\_rank\_r\_sanity.py},
\texttt{day10\_general\_ht.py}, \texttt{day11\_task\_dep\_Dt.py},
\texttt{day12\_achievability.py}, and
\texttt{day13\_achievability\_fixes.py}.
\end{abstract}

\section{Setting}\label{sec:setting}

Fix integers $T \geq 1$, $d \geq 2$, $1 \leq r \leq d$, a radius
$B > 0$, and a rate $R \geq 0$ (bits).

\paragraph{Task vectors and task metrics.} Each task is a pair
$(\tau_t, H_t) \in \R^d \times \R^{d\times d}$ with $H_t \succeq 0$ and
$\rank(H_t) \leq r$. Let $V_t := \range(H_t) \subseteq \R^d$ be the
task-$t$ subspace (so $\dim V_t \leq r$). We assume the affine
normalization $\tau_t \in V_t$, i.e.\ $\tau_t = P_{V_t}\tau_t$, and
the scale normalization $\tau_t^\top H_t \tau_t \leq B^2$.\footnote{%
  The scale normalization is a $B$-in-the-$H_t$-metric bound;
  for the Frobenius-norm-bounded LoRA special case $H_t = P_{V_t}$
  it reduces to $\|\tau_t\|^2 \leq B^2$ since
  $\tau_t\in V_t$.}

\paragraph{Per-task loss and max-distortion.} Per-task loss
$f_t(w) = \tfrac12 (w-\tau_t)^\top H_t (w-\tau_t)$ attains its minimum
zero at any $w$ with $P_{V_t} w = \tau_t$. The \emph{max-distortion}
figure of merit is
\[
  \Delta(w^\star; \boldsymbol\tau, \boldsymbol H)
  \;=\; \max_{t\in[T]}\; (w^\star - \tau_t)^\top H_t (w^\star - \tau_t).
\]

\paragraph{Merging algorithms and RD function.} A rate-$R$ merging
algorithm is $(E,D)$ with $E:(\R^d)^T\to\{1,\dots,2^R\}$ and
$D:\{1,\dots,2^R\}\to\R^d$, producing $w^\star = D(E(\boldsymbol\tau))$.
The rank-$r$ rate--distortion function is
\begin{equation}\label{eq:rd-def-rankr}
  \mathcal{D}^\star(T, d, r, B, R)
  \;=\; \inf_{E,D}\;\sup_{(\boldsymbol\tau, \boldsymbol H)}\;
  \E\bigl[\Delta(w^\star;\boldsymbol\tau,\boldsymbol H)\bigr],
\end{equation}
where the sup is over admissible $(\boldsymbol\tau, \boldsymbol H)$
in the sense above.

\paragraph{LoRA special case.} For a LoRA fine-tune with rank-$r$
update $\tau_t = A_t B_t^\top$ (here $A_t, B_t \in \R^{d_0\times r}$
with $d = d_0^2$), weight-disentanglement gives $H_t = P_{V_t}$
with $V_t = \mathrm{span}(\text{columns of }A_t)\otimes\mathrm{span}
(\text{columns of }B_t)$ (Ortiz-Jimenez et~al., NeurIPS~2023, Def.\ 3.1;
see \texttt{notes/ortiz\_jimenez\_tangent.md}~\S4). For this scaffold
we will specialize to $H_t = P_{V_t}$ whenever a concrete computation
is needed; the general $H_t \succeq 0$ case is an extension.

\section{$H_t$-weighted max $\geq$ avg identity}\label{sec:id}

The Phase 0 proof (Theorem~1 of \texttt{toy\_theorem\_v0.tex}) used
the deterministic identity
$T^{-1}\sum_t\|w - \tau_t\|^2 = \|w - \bar\tau\|^2 + T^{-1}\sum_t
\|\tau_t - \bar\tau\|^2$, provable because $\sum_t(\tau_t-\bar\tau) = 0$
identically. We generalize to $H_t$-weighted distortion.

\paragraph{The $H_t$-weighted centroid.} Let
$\Hbar := T^{-1}\sum_t H_t \succeq 0$ and define
\begin{equation}\label{eq:htau}
  \tauH \;:=\; \Hbar\hpinv \cdot T^{-1}\sum_{t=1}^T H_t \tau_t,
\end{equation}
where $(\cdot)\hpinv$ denotes the Moore--Penrose pseudoinverse.
Equivalently, $\tauH$ is the unique element of
$\range(\Hbar) = V_1 + \dots + V_T$ satisfying $\Hbar \tauH =
T^{-1}\sum_t H_t\tau_t$. This is the $H_t$-weighted analog of
$\bar\tau = T^{-1}\sum_t\tau_t$; the two agree when $H_t = I_d$ for
all $t$.

\begin{lemma}[$H_t$-weighted max $\geq$ avg identity]\label{lem:id}
For any $w \in \R^d$ and any admissible $(\boldsymbol\tau,
\boldsymbol H)$,
\begin{equation}\label{eq:id}
  \tfrac{1}{T}\sum_{t=1}^T (w-\tau_t)^\top H_t (w-\tau_t)
  \;=\; (w - \tauH)^\top\Hbar(w-\tauH)
  \;+\; \tfrac{1}{T}\sum_{t=1}^T (\tau_t-\tauH)^\top H_t (\tau_t-\tauH).
\end{equation}
Hence
\begin{equation}\label{eq:maxavg}
  \max_t (w-\tau_t)^\top H_t (w-\tau_t)
  \;\geq\; (w - \tauH)^\top\Hbar(w-\tauH)
  \;+\; \tfrac{1}{T}\sum_t (\tau_t-\tauH)^\top H_t (\tau_t-\tauH).
\end{equation}
\end{lemma}

\begin{proof}
Expand the LHS of \eqref{eq:id}:
\[
  \tfrac1T\sum_t (w-\tau_t)^\top H_t (w-\tau_t)
  \;=\; w^\top \Hbar w - 2 w^\top\!\Bigl(\tfrac1T\sum_t H_t\tau_t\Bigr)
    + \tfrac1T\sum_t \tau_t^\top H_t \tau_t.
\]
Using $\tfrac1T\sum_t H_t\tau_t = \Hbar\tauH$ (by
\eqref{eq:htau} and $\Hbar\Hbar\hpinv|_{\range(\Hbar)} = I$),
\[
  = w^\top\Hbar w - 2 w^\top\Hbar\tauH + \tfrac1T\sum_t\tau_t^\top H_t \tau_t
  \;=\; (w-\tauH)^\top\Hbar(w-\tauH) - \tauH^\top\Hbar\tauH +
    \tfrac1T\sum_t\tau_t^\top H_t\tau_t.
\]
The last two terms equal $T^{-1}\sum_t(\tau_t - \tauH)^\top H_t
(\tau_t-\tauH)$, by the same expansion applied with $w\leftarrow\tauH$.
\eqref{eq:maxavg} follows from max $\geq$ avg.
\end{proof}

\begin{remark}[Why the cross-term vanishes]
The Phase 0 cross-term $\sum_t(\tau_t-\bar\tau) = 0$ is the
\emph{linear} cross-term. Its $H_t$-weighted analog is
$\sum_t H_t(\tau_t-\tauH) = T(\Hbar\tauH - \Hbar\tauH) = 0$, which is
zero by the \emph{definition} \eqref{eq:htau} of $\tauH$. No
probabilistic step is needed. At $H_t = I_d$ we recover the Phase 0
identity exactly.
\end{remark}

\begin{remark}[Components of $w$ outside $\range(\Hbar)$]
\label{rem:project}
Any component of $w$ orthogonal to $\range(\Hbar) = V_1 + \dots + V_T$
annihilates against every $H_t$ and does not affect
\eqref{eq:id}--\eqref{eq:maxavg}. Without loss of generality, the
decoder may output $w^\star \in V_1 + \dots + V_T$, which has
dimension at most $\deff := \rank(\sum_t P_{V_t}) \leq \min(Tr, d)$.
This is the key dimensional saving behind the rank-$r$ rate--distortion
improvement.
\end{remark}

\section{Hard distribution}\label{sec:hard}

For the lower bound we pick a hard distribution $P$ on
$(\boldsymbol\tau, \boldsymbol H)$ tuples.

\paragraph{Stiefel-random LoRA subspaces.} Independently for each
$t\in[T]$:
\begin{enumerate}
  \item Sample $U_t \in \Stiefel(d, r)$ uniformly (columns an
  orthonormal basis of $V_t$; equivalently, QR of a $d\times r$
  iid Gaussian).
  \item Set $H_t = P_{V_t} = U_t U_t^\top$.
  \item Sample $v_t \sim \mathrm{Unif}(B \cdot S^{r-1})$ inside
  $\R^r$; set $\tau_t = U_t v_t \in V_t$.
\end{enumerate}
Under $P$, each $\tau_t$ has $\|\tau_t\|^2 = B^2$, lives in its own
random $r$-dimensional subspace, and the pairwise subspace overlap
scales as $r^2/d$ in expectation.

\paragraph{Overlap coefficient.} Define
\begin{equation}\label{eq:gamma}
  \gamma := \bigl\|T^{-1}\textstyle\sum_t P_{V_t}\bigr\|_{\mathrm{op}}
  \;\in\; [1/T,\, 1].
\end{equation}
$\gamma = 1/T$ iff the $V_t$'s are mutually orthogonal (rank of sum
$= Tr$); $\gamma = 1$ iff all $V_t = V$ coincide (rank of sum $= r$).
For Stiefel-random $V_t$ with $Tr \leq d$, $\gamma \approx r/d +
o(1)$ typically. As Lemma~\ref{lem:floor-rankr} will show, the scalar
that actually controls the lower bound is $\deff$ itself;
$\gamma$ enters only through $\deff$'s dependence on the overlap
structure.

\section{Closed-form floor under $P$}\label{sec:floor}

The expected floor
$\E_P[T^{-1}\sum_t(\tau_t-\tauH)^\top P_{V_t}(\tau_t-\tauH)]$ from
Lemma~\ref{lem:id} admits an exact closed form in terms of $\deff$
alone.

\begin{lemma}[Closed-form floor]\label{lem:floor-rankr}
Under the hard distribution $P$ of \S\ref{sec:hard} with
$H_t = P_{V_t}$,
\begin{equation}\label{eq:floor-closed}
  \E_P\!\left[\tfrac{1}{T}\sum_t (\tau_t - \tauH)^\top P_{V_t}(\tau_t - \tauH)\right]
  \;=\; B^2\!\left(1 - \frac{\E_P[\deff]}{Tr}\right),
\end{equation}
where $\deff := \rank\bigl(\sum_t P_{V_t}\bigr) \leq \min(Tr, d)$.
\end{lemma}

\begin{proof}
Two steps: a deterministic identity reducing the floor to a bilinear
form in $\bar\tau$, then an expectation-via-isotropy computation.

\emph{Step 1 (deterministic identity).} Apply Lemma~\ref{lem:id} at
$w = \tauH$; the centroid term vanishes. Using $P_{V_t}\tau_t = \tau_t$
(since $\tau_t\in V_t$), $T^{-1}\sum_t P_{V_t} = \bar H$, and
$T^{-1}\sum_t\tau_t =: \bar\tau$:
\[
  \tfrac1T\sum_t(\tau_t - \tauH)^\top P_{V_t}(\tau_t-\tauH)
  \;=\; B^2 - 2\bar\tau^\top\tauH + \tauH^\top\bar H\tauH.
\]
By definition \eqref{eq:htau} combined with $P_{V_t}\tau_t = \tau_t$,
$\bar H\tauH = T^{-1}\sum_t P_{V_t}\tau_t = \bar\tau$. Hence
$\bar\tau^\top\tauH = \tauH^\top\bar H\tauH$, and since $\tauH = \bar H^+
\bar\tau$ on $\range(\bar H)$, $\tauH^\top\bar H\tauH = \bar\tau^\top
\bar H^+\bar\tau$. Substituting,
\begin{equation}\label{eq:floor-det}
  \tfrac1T\sum_t(\tau_t-\tauH)^\top P_{V_t}(\tau_t-\tauH)
  \;=\; B^2 - \bar\tau^\top\bar H^+\bar\tau,
\end{equation}
which is a deterministic identity on admissible tuples.

\emph{Step 2 (expectation via isotropy).} Condition on $V_1,\dots,V_T$.
Under $P$, $\tau_t | V_t$ is uniform on $B\cdot S^{r-1}\subset V_t$
independently across $t$, so $\E[\tau_t | V_t] = 0$ and
$\E[\tau_t\tau_t^\top | V_t] = (B^2/r)\,P_{V_t}$. Expanding
$\bar\tau^\top\bar H^+\bar\tau = T^{-2}\sum_{s,t}\tau_s^\top
\bar H^+\tau_t$: off-diagonal ($s\neq t$) terms vanish in expectation
by independence and zero conditional mean; the diagonal gives
\[
  \E[\tau_t^\top\bar H^+\tau_t \mid V]
  \;=\; \tr\!\bigl(\bar H^+\cdot(B^2/r)P_{V_t}\bigr)
  \;=\; (B^2/r)\,\tr(\bar H^+ P_{V_t}).
\]
Summing over $t$ and using $\sum_t P_{V_t} = T\bar H$,
\[
  \E[\bar\tau^\top\bar H^+\bar\tau \mid V]
  \;=\; \frac{B^2}{T^2 r}\,\tr(\bar H^+\cdot T\bar H)
  \;=\; \frac{B^2}{Tr}\,\tr(\bar H^+\bar H)
  \;=\; \frac{B^2\deff}{Tr},
\]
using $\bar H^+\bar H = P_{\range(\bar H)}$ with trace $\deff$. Taking
expectation over $V$ and substituting into \eqref{eq:floor-det}
yields \eqref{eq:floor-closed}.
\end{proof}

\begin{remark}[Endpoint consistency]\label{rem:floor-endpoints}
Specializing $\deff$ to the four regimes tracked by \S\ref{sec:limits}
reproduces each sanity-check proposition verbatim:
\begin{itemize}
  \item $r = d$: $V_t = \R^d$, $\deff = d$, $Tr = Td$; floor
    $= B^2(1-1/T)$ (Prop.~\ref{prop:full}).
  \item $T = 1$: $\deff = r$, $Tr = r$; floor $= 0$
    (Prop.~\ref{prop:T1}).
  \item $V_t = V$ shared: $\deff = r$; floor $= B^2(1-1/T)$
    (Prop.~\ref{prop:shared}).
  \item \emph{Stiefel-random with $Tr \leq d$}: a.s.\ $\deff = Tr$
    (the subspaces are in general position), so $f_{\mathrm{floor}} = 0$.
    Geometrically, each $\tau_t$ is exactly reconstructible from
    $\tauH$ through its own projector $P_{V_t}$; the multi-task floor
    vanishes in the generic-position regime.
\end{itemize}
The Day-8 open-question worry that the orthogonal-limit floor might
collapse to $B^2$ is resolved in the opposite direction: it collapses
to $0$, and \eqref{eq:floor-closed} interpolates monotonically via
$\deff$.
\end{remark}

\section{Main theorem: rank-$r$ RD lower bound}\label{sec:thm}

We promote Conjecture~2 (Day-8 scaffold) to a theorem for the
LoRA-consistent case $H_t = P_{V_t}$, leveraging
Lemma~\ref{lem:floor-rankr} for the floor and a rank-$r$ analog of
Phase 0 Lemma~2 for the compression term.

\begin{theorem}[Rank-$r$ RD lower bound, $H_t = P_{V_t}$]\label{thm:rankr}
Let $T\geq 1$, $d\geq 2$, $1\leq r\leq d$, $B > 0$, $R\geq 0$. Fix a
distribution on $(V_1,\dots,V_T)$ under which $\deff := \rank(\sum_t
P_{V_t})$ is a.s.\ constant, and sample $\tau_t = U_t v_t$ with
$v_t \sim \mathrm{Unif}(B\cdot S^{r-1})$ conditional on $V_t$. Then
for the admissible class of \S\ref{sec:setting} with $H_t = P_{V_t}$,
\begin{equation}\label{eq:thm}
  \mathcal{D}^\star(T, d, r, B, R)
  \;\geq\;
  B^2\!\left(1 - \frac{\deff}{Tr}\right)
  \;+\; c_{\mathrm{TQ}}\cdot\frac{B^2\deff}{Tr}
  \cdot 2^{-2R/\deff}.
\end{equation}
Under the Stiefel-random hard distribution of \S\ref{sec:hard} with
$Tr \leq d$, $\deff = Tr$ a.s., so
\begin{equation}\label{eq:thm-stiefel}
  \mathcal{D}^\star(T, d, r, B, R)
  \;\geq\;
  c_{\mathrm{TQ}}\cdot B^2\cdot 2^{-2R/(Tr)}.
\end{equation}
The constant $c_{\mathrm{TQ}}$ is the universal constant of
\textsc{TurboQuant} Theorem~3 (arXiv~2504.19874, \S3.3), same as in
Theorem~1 of \texttt{theory/toy\_theorem\_v0.tex}.
\end{theorem}

\begin{proof}
The proof mirrors Phase 0 in six steps: Yao's minimax, the max-$\geq$-avg
reduction via Lemma~\ref{lem:id}, restriction of the decoder to
$\range(\bar H)$, change of variable to Euclidean geometry in that
subspace, a rank-$r$ analog of Phase 0 Lemma~2, and combination.

\emph{Step 1 (Yao).} For any $P$ supported on the admissible class,
$\mathcal{D}^\star \geq \inf_{E,D}\E_P[\max_t(w^\star - \tau_t)^\top
P_{V_t}(w^\star - \tau_t)]$. Take $P$ as in the hypothesis.

\emph{Step 2 (max $\geq$ avg via Lemma~\ref{lem:id}).}
\eqref{eq:maxavg} with $H_t = P_{V_t}$ is deterministic in the tuple.
Taking $\E_P$ and using Lemma~\ref{lem:floor-rankr} on the floor term,
\begin{equation}\label{eq:reduce-rankr}
  \E_P\!\left[\max_t(w^\star-\tau_t)^\top P_{V_t}(w^\star-\tau_t)\right]
  \;\geq\;
  \E_P\!\left[(w^\star-\tauH)^\top\bar H(w^\star-\tauH)\right]
  + B^2\!\left(1 - \frac{\deff}{Tr}\right).
\end{equation}

\emph{Step 3 (restrict decoder to $\range(\bar H)$).} By
Remark~\ref{rem:project}, components of $w^\star$ orthogonal to
$\range(\bar H)$ don't affect any $H_t$-weighted distortion, so
without loss of generality the decoder emits
$w^\star \in \range(\bar H)$.

\emph{Step 4 (change of variable).} Let $\xi := \bar H^{1/2}(w^\star
- \tauH)$ and $\eta := \bar H^{1/2}\tauH$, both in
$\range(\bar H)\cong\R^{\deff}$. Then $(w^\star-\tauH)^\top
\bar H(w^\star-\tauH) = \|\xi\|^2 = \|\bar H^{1/2}w^\star - \eta\|^2$.
$\bar H^{1/2}$ on $\range(\bar H)$ is a bijection; $w^\star$ is a
rate-$R$ function of the tuple, so $\bar H^{1/2}w^\star$ is too. The
reduced problem is to quantize $\eta$ at rate $R$ minimizing
$\E\|\bar H^{1/2}w^\star-\eta\|^2$ in $\R^{\deff}$.

\emph{Step 5 (rank-$r$ Lemma~2 analog).} Condition on
$W := \range(\bar H) = V_1 + \cdots + V_T$, a $\deff$-dim subspace of
$\R^d$. By the $O(d)$-invariance of the Stiefel-random sampling, the
joint conditional law of $(V_1,\dots,V_T,\boldsymbol\tau)$ given
$\{V_1+\cdots+V_T = W\}$ is $O(W)$-invariant: for any $\tilde O\in O(W)$
extended to $O\in O(d)$ by $O|_{W^\perp} = I$, the pushforward
$(OV_t, O\tau_t)_t$ has the same conditional law (the $V_t$'s are
permuted within $W$ by $O$, but their joint law is unchanged because
the unconditional Stiefel measure is $O(d)$-invariant and $W$ is
preserved setwise; the $v_t$'s sit in the uniform-sphere law of their
respective $V_t$'s which is $O$-equivariant). Under this action
$\bar H \mapsto O\bar H O^\top$, $\bar H^{1/2}\mapsto O\bar H^{1/2} O^\top$,
$\tauH\mapsto O\tauH$, so $\eta \mapsto O\eta$. Hence
$\eta\,|\,W \stackrel{d}{=} \tilde O\,\eta\,|\,W$ for every
$\tilde O\in O(W)$, i.e.\ the conditional law of $\eta$ given $W$ is
$O(W)$-rotationally invariant in $W \cong \R^{\deff}$.

Now apply Lemma~2 of \texttt{theory/toy\_theorem\_v0.tex}
(scale-invariant form, equation \texttt{eq:rd-rot-scale} there) to
the rate-$R$ reconstruction $\bar H^{1/2}w^\star$ of the
$O(W)$-invariant source $\eta$ in the $\deff$-dim subproblem:
\[
  \E\!\left[\|\bar H^{1/2}w^\star - \eta\|^2 \,\big|\, W\right]
  \;\geq\;
  c_{\mathrm{TQ}}\cdot\E\!\left[\|\eta\|^2\,\big|\,W\right]
  \cdot 2^{-2R/\deff}\cdot(1 - o_{\deff}(1)),
\]
where the $(1-o_{\deff}(1))$ prefactor is the Phase 0 concentration
correction (Remark~3 of \texttt{theory/toy\_theorem\_v0.tex}) and is
absorbed into $c_{\mathrm{TQ}}$ by convention. Empirical verification
of the $O(W)$-invariance claim is in
\texttt{code/synthetic/day9\_verify\_lemma6.py}: the conditional
second-moment matrix $\E[\eta_W\eta_W^\top|W]$ is scalar within
Marchenko--Pastur sampling bounds at $n=10^5$, and
$\E[\|\eta\|^2|W] = B^2\deff/(Tr)$ is matched to four digits.

\emph{Step 6 (combine).} Take expectation over $W$ (and hence over the
underlying $V$). By Lemma~\ref{lem:floor-rankr},
\[
  \E_P[\|\eta\|^2]
  \;=\; \E_P[\tauH^\top\bar H\tauH]
  \;=\; B^2 - B^2\!\left(1-\frac{\deff}{Tr}\right)
  \;=\; \frac{B^2\deff}{Tr}.
\]
Substituting into Step 5 and combining with \eqref{eq:reduce-rankr}
yields \eqref{eq:thm}. Under Stiefel-random sampling with $Tr\leq d$,
$\deff = Tr$ a.s., and \eqref{eq:thm-stiefel} follows.
\end{proof}

\begin{remark}[Interpretation]
The compression prefactor $B^2\deff/(Tr)$ reads as the per-task
``compressible budget'': of the total second moment $B^2$ per task,
$B^2\deff/(Tr)$ is centroid-captured (available for rate to act on)
and $B^2(1-\deff/(Tr))$ is locked into the floor. At $V_t=V$:
budget $B^2/T$, floor $B^2(1-1/T)$ — recovers Phase 0. At Stiefel-
random with $Tr\leq d$: budget $B^2$, floor $0$ — the generic-position
regime has no floor and all rate goes against the full per-task
variance.
\end{remark}

\begin{remark}[The $Tr > d$ regime]
Theorem~\ref{thm:rankr} states a bound parameterized by an a.s.\
constant $\deff$. For Stiefel-random sampling with $Tr > d$, $\deff =
d$ a.s.\ (the subspaces jointly span $\R^d$), and \eqref{eq:thm}
yields floor $B^2(1-d/(Tr)) > 0$ and compression term
$c_{\mathrm{TQ}} (B^2 d/(Tr)) 2^{-2R/d}$. The LoRA-relevant regime is
$r\ll d$, so $Tr\leq d$ comfortably holds.
\end{remark}

\section{Extension to general $H_t \succeq 0$}\label{sec:general-ht}

Theorem~\ref{thm:rankr} uses $H_t = P_{V_t}$. For the
$\text{MSE}\to\text{CE}$ bridge the target is general $H_t \succeq 0$
with $\range(H_t)\subseteq V_t$ (LoRA Fisher information is rank-$r$
but not a projector; its spectrum inside $V_t$ is the task-$t$
curvature, and may differ across tasks). We show that under the
natural generalization of the hard distribution, the bound of
Theorem~\ref{thm:rankr} holds verbatim, for \emph{arbitrary
per-task} eigenvalue spectra $(D_t)_{t=1}^T$.

\paragraph{Generalized hard distribution $P^\star$.} Fix
\emph{per-task} $r\times r$ positive-definite matrices $D_1, \dots,
D_T$ with eigenvalues in $(0, +\infty)$.\footnote{%
  Each $D_t$ encodes the task-$t$ curvature spectrum in a canonical
  coordinate system of $V_t$. The choice $D_t = I_r$ for all $t$
  recovers the hard distribution $P$ of \S\ref{sec:hard}; the choice
  $D_t \equiv D_0$ common across tasks is the special case
  (``common-spectrum'') highlighted on Day 10 of the scaffold
  development. Nothing in the LB below uses any relationship among
  the $D_t$'s.} Independently for each $t\in[T]$:
\begin{enumerate}
  \item Sample $U_t \in \Stiefel(d, r)$ uniformly (as in \S\ref{sec:hard}).
  \item Set $H_t = U_t D_t U_t^\top$, so $\range(H_t) = V_t =
        \mathrm{col}(U_t)$ and $\rank(H_t) = r$.
  \item Sample $u_t \sim \mathrm{Unif}(S^{r-1})$ and set $v_t =
        B\,D_t^{-1/2} u_t \in \R^r$; set $\tau_t = U_t v_t\in V_t$.
\end{enumerate}
By construction $\tau_t^\top H_t \tau_t = v_t^\top D_t v_t = B^2$
per task (saturates the scale normalization of
\S\ref{sec:setting}), and $\tau_t$ is uniform on the $H_t$-ellipsoid
$\{\tau\in V_t : \tau^\top H_t\tau = B^2\}$. For $D_t \equiv I_r$,
$P^\star$ reduces to $P$.

\begin{lemma}[Floor formula under $P^\star$]\label{lem:floor-general}
Under $P^\star$ with any fixed positive-definite matrices
$(D_t)_{t=1}^T$ (possibly distinct across tasks),
\[
  \E_{P^\star}\!\left[\tfrac1T\sum_t(\tau_t-\tauH)^\top H_t(\tau_t-\tauH)\right]
  \;=\; B^2\!\left(1 - \frac{\E[\deff]}{Tr}\right),
\]
with $\deff = \rank(\sum_t P_{V_t})$ as in Lemma~\ref{lem:floor-rankr}.
\end{lemma}

\begin{proof}
The Lemma~\ref{lem:floor-rankr} proof has two steps; only Step~2
(isotropy) needs checking.

\emph{Step 1 (same as Lemma~\ref{lem:floor-rankr}).} Apply
Lemma~\ref{lem:id} at $w=\tauH$; using $\tau_t^\top H_t\tau_t = B^2$
and $\mu := T^{-1}\sum_t H_t\tau_t = \bar H\tauH$,
\begin{equation}\label{eq:floor-general-det}
  \tfrac1T\sum_t(\tau_t-\tauH)^\top H_t(\tau_t-\tauH)
  \;=\; B^2 - \tauH^\top\bar H\tauH
  \;=\; B^2 - \mu^\top\bar H^+\mu,
\end{equation}
using $\tauH^\top\bar H\tauH = \mu^\top\bar H^+\bar H\bar H^+\mu =
\mu^\top\bar H^+\mu$ on $\range(\bar H)$.

\emph{Step 2 (isotropy under $P^\star$).} Conditional on $U_t$ (hence
on $V_t$ and on $H_t$), $v_t = B D_t^{-1/2}u_t$ with $u_t$ uniform on
$S^{r-1}$; hence
\[
  \E[v_t v_t^\top \mid U_t] = B^2\,D_t^{-1/2}\,\E[u_t u_t^\top]\,D_t^{-1/2}
  \;=\; (B^2/r)\,D_t^{-1}.
\]
Therefore
$\E[\tau_t\tau_t^\top\mid U_t] = U_t\,\E[v_t v_t^\top]\,U_t^\top =
(B^2/r)\,U_t D_t^{-1} U_t^\top = (B^2/r)\,H_t^+$ (using
$H_t^+ = U_t D_t^{-1}U_t^\top$ on $V_t$). Then
\[
  \E[H_t\tau_t\tau_t^\top H_t\mid U_t]
  \;=\; (B^2/r)\,H_t\,H_t^+\,H_t \;=\; (B^2/r)\,H_t,
\]
by the PSD identity $H_t H_t^+ H_t = H_t$, which holds per-task
regardless of $D_t$. Summing over $t$ and using
off-diagonal terms vanish (independence, zero conditional mean),
\[
  \E[\mu\mu^\top\mid V]
  \;=\; T^{-2}\sum_t\E[H_t\tau_t\tau_t^\top H_t\mid U_t]
  \;=\; (B^2/(T^2 r))\sum_t H_t
  \;=\; (B^2/(Tr))\,\bar H.
\]
Hence $\E[\mu^\top\bar H^+\mu \mid V] = (B^2/(Tr))\tr(\bar H^+\bar H)
= B^2\deff/(Tr)$. Take expectation over $V$ and substitute into
\eqref{eq:floor-general-det}.
\end{proof}

\begin{theorem}[Rank-$r$ RD lower bound, general $H_t$]\label{thm:general}
Under $P^\star$ with any fixed positive-definite matrices
$(D_t)_{t=1}^T$ (possibly distinct across tasks), and with
$\rank(H_t) = r$ for all $t$,
\begin{equation}\label{eq:thm-general}
  \mathcal{D}^\star(T,d,r,B,R)
  \;\geq\;
  B^2\!\left(1-\frac{\deff}{Tr}\right)
  \;+\; c_{\mathrm{TQ}}\cdot\frac{B^2\deff}{Tr}\cdot 2^{-2R/\deff}.
\end{equation}
In particular, for Stiefel-random $U_t$ with $Tr\leq d$, $\deff = Tr$
a.s.\ and the floor vanishes, giving
$\mathcal{D}^\star \geq c_{\mathrm{TQ}}\cdot B^2\cdot 2^{-2R/(Tr)}$.
\end{theorem}

\begin{proof}
Six steps exactly as in Theorem~\ref{thm:rankr}; we only note the
modifications.

\emph{Steps 1--4} are identical: Yao~$\to$~Lemma~\ref{lem:id}~$\to$~
restriction to $\range(\bar H)$~$\to$~change of variable
$\eta:=\bar H^{1/2}\tauH$. After Step~4 the sub-problem is rate-$R$
quantization of $\eta$ in $\R^{\deff}$ under squared Euclidean loss,
with Step~2's floor term now given by Lemma~\ref{lem:floor-general}.

\emph{Step 5 (rotational invariance under $P^\star$).} Condition on
$W := V_1 + \cdots + V_T$, a $\deff$-dim subspace of $\R^d$. For
$\tilde O\in O(W)$ extended to $O\in O(d)$ by $O|_{W^\perp} = I$, the
tuple $(U_t, v_t)_t$ transforms as follows under the pushforward:
$U_t \mapsto OU_t$ (again Stiefel-distributed, since the Stiefel
measure is $O(d)$-invariant and $W$ is preserved); $v_t$ is unchanged
(it lives in $\R^r$, not $\R^d$, and its law is $O(d)$-free, noting
that $v_t = B D_t^{-1/2} u_t$ is determined by $u_t$ and the
task-$t$-labeled $D_t$, neither of which transforms under $O$).
Hence $H_t = U_t D_t U_t^\top \mapsto (OU_t) D_t (OU_t)^\top =
O H_t O^\top$, $\tau_t = U_t v_t \mapsto O\tau_t$, $\bar H \mapsto
O\bar H O^\top$, $\bar H^{1/2} \mapsto O\bar H^{1/2}O^\top$, $\mu
\mapsto O\mu$, $\tauH \mapsto O\tauH$, $\eta \mapsto O\eta$. The
joint conditional law of the tuple given $W$ is therefore
$O(W)$-invariant (regardless of whether the $D_t$'s match across
tasks), and so is the marginal of $\eta$ given $W$. Apply Lemma~2
of \texttt{toy\_theorem\_v0.tex} as in Theorem~\ref{thm:rankr}'s
Step~5; use $\E[\|\eta\|^2\mid W] = B^2\deff/(Tr)$ from
Lemma~\ref{lem:floor-general}.

\emph{Step 6 (combine).} Identical to Theorem~\ref{thm:rankr}'s
Step~6; yields \eqref{eq:thm-general}.
\end{proof}

\begin{remark}[Why the bound doesn't depend on $(D_t)_t$]
For the LB the only role of $D_t$ is in the distribution of $\tau_t$:
the per-task $D_t^{-1/2}$ stretch ensures $\E[H_t\tau_t\tau_t^\top
H_t \mid U_t] = (B^2/r)H_t$ regardless of the eigenvalue spectrum,
so the second moment of $\mu$ inside $\range(\bar H)$ is always
$(B^2/(Tr))\bar H$. Geometrically, sampling $\tau_t$ uniformly on
the $H_t$-ellipsoid makes the ``$H_t$-weighted energy'' isotropic in
$V_t$, and this ``per-task isotropy'' is what the LB needs --- not
any relationship among the $D_t$'s. An alternative distribution
(e.g.\ $\tau_t$ uniform on the Euclidean sphere $B\cdot S^{r-1}\cap
V_t$) breaks this isotropy when $D_t \neq I$, and the floor formula
changes. Numerical checks across seven $D_0$/$D_t$ configurations
(common, per-task-mismatched, four spectrum shapes): Day-10
\texttt{day10\_general\_ht.py} (common $D_0$) and Day-11
\texttt{day11\_task\_dep\_Dt.py} (task-dependent $D_t$) both agree
with Lemma~\ref{lem:floor-general} to four to five digits and
respect the Marchenko--Pastur rotational-invariance bound.
\end{remark}

\begin{remark}[Practical $H_t$: CE-loss Fisher]
For LLM cross-entropy loss, $H_t$ near $\tau_t$ is the Fisher
information matrix, which is rank-$r$ PSD with $\range(H_t)
\subseteq V_t$ and (generally) a task-dependent eigenvalue spectrum
inside $V_t$. Theorem~\ref{thm:general} applies directly to this
setting: the proof uses only per-task properties of $D_t$
(specifically the identity $H_t H_t^+ H_t = H_t$ and the
$O(d)$-freeness of $v_t$), never any relationship across tasks.
This closes the $\text{MSE}\to\text{CE}$ bridge under the
$P^\star$ hard distribution for LoRA Fisher metrics with
heterogeneous task curvature.
\end{remark}

\section{Sanity-check corollaries}\label{sec:limits}

The three Day-8 sanity-check limits are now corollaries of
Theorem~\ref{thm:rankr} + Lemma~\ref{lem:floor-rankr}. Each reduction
is a one-line substitution of $\deff$.

\begin{proposition}[Full-rank limit $r=d$, $V_t = \R^d$]\label{prop:full}
At $r=d$: $V_t = \R^d$, $H_t = I_d$, $P$ collapses to
$\mathrm{Unif}(B\cdot S^{d-1})^{\otimes T}$, $\tauH = \bar\tau$,
$\deff = d$ a.s. Theorem~\ref{thm:rankr} with $\deff = d$ and
$Tr = Td$ yields
\[
  \mathcal{D}^\star(T,d,d,B,R)
  \;\geq\; B^2(1-1/T) + c_{\mathrm{TQ}}(B^2/T) 2^{-2R/d},
\]
which is Theorem~1 of \texttt{theory/toy\_theorem\_v0.tex} verbatim
(equation \texttt{eq:thm-main} there).
\end{proposition}

\begin{proposition}[Single-task limit $T=1$]\label{prop:T1}
At $T=1$: $\bar H = P_{V_1}$, $\tauH = \tau_1$ (since $\tau_1\in V_1$),
floor $= 0$, $\deff = r$. Theorem~\ref{thm:rankr} with $\deff = r$
and $Tr = r$ yields
\[
  \mathcal{D}^\star(1, d, r, B, R) \;\geq\; c_{\mathrm{TQ}}\cdot B^2
  \cdot 2^{-2R/r},
\]
which is \textsc{TurboQuant} Theorem~3 applied to
$\mathrm{Unif}(B\cdot S^{r-1})\subset V_1$ with identical constant.
\end{proposition}

\begin{proposition}[Shared-subspace limit $V_t = V$]\label{prop:shared}
If $V_t = V$ deterministically with $\dim V = r$: $\bar H = P_V$,
$\tauH = \bar\tau\in V$, $\deff = r$ a.s. Theorem~\ref{thm:rankr} with
$\deff = r$ yields
\[
  \mathcal{D}^\star_{\text{shared}}(T, r, B, R)
  \;\geq\; B^2(1-1/T) + c_{\mathrm{TQ}}(B^2/T) 2^{-2R/r},
\]
which is Theorem~1 of \texttt{theory/toy\_theorem\_v0.tex} with
$d \mapsto r$.
\end{proposition}

\paragraph{Summary table.}
\begin{center}
\begin{tabular}{lcccc}
\hline
Regime                         & $\deff$ & Floor           & Compression prefactor & Decay exponent \\
\hline
$r = d$                        & $d$     & $B^2(1-1/T)$    & $B^2/T$               & $2^{-2R/d}$    \\
$T = 1$                        & $r$     & $0$             & $B^2$                 & $2^{-2R/r}$    \\
$V_t = V$ shared               & $r$     & $B^2(1-1/T)$    & $B^2/T$               & $2^{-2R/r}$    \\
Stiefel-random, $Tr \leq d$    & $Tr$    & $0$             & $B^2$                 & $2^{-2R/(Tr)}$ \\
\hline
\end{tabular}
\end{center}
In every row floor plus compression-prefactor sum to $B^2$, reflecting
the conservation $f_{\mathrm{floor}}\cdot B^2 + \E[\|\eta\|^2] = B^2$
from the proof of Theorem~\ref{thm:rankr}.

\section{Numerical evidence}\label{sec:num}

Numerics in \texttt{code/synthetic/day8\_rank\_r\_sanity.py} sweep
$(T, d, r, b, \text{overlap}) \in \{2,4\} \times \{128, 512\}
\times \{4, 8, 16, 32, d\} \times \{1, 2, 4, 8\} \times
\{\text{random}, \text{shared}\}$. Three checks, all tied to
Theorem~\ref{thm:rankr} and Lemma~\ref{lem:floor-rankr}:
\begin{itemize}
  \item \textbf{(a) Regression at $r = d$} (Prop.~\ref{prop:full}):
  empirical Chebyshev floor divided by $B^2(1-1/T)$ should lie in
  $[0.998, 1.017]$, the range observed in Phase 0
  (\texttt{day5\_lower\_bound\_sanity.py}). Passed Day~8.
  \item \textbf{(b) Shared-subspace slope} (Prop.~\ref{prop:shared}):
  on the quantized-$\tauH$ excess above the empirical floor, the
  $\log_2\text{excess}$ vs.\ $b$ slope should equal $-2$
  per-bit-per-coord in the $\deff = r$ coordinate system, i.e.\
  $-2 d/r$ in the ambient $d$-coord basis after the natural
  subspace-aware quantizer. Day~8 uniform-$[-B,B]$ quantization
  saturated at low bit-rate (per-coord std $B/\sqrt{Tr} \ll B$);
  Day~9 replaces it with scale-adaptive quantization matched to
  Lemma~\ref{lem:floor-rankr}'s per-coord scale $\sigma_{\mathrm{pc}}
  = B/\sqrt{Tr}$.
  \item \textbf{(c) Random-subspace $\deff$}: empirical $\deff$ from
  the slope fit should match $\rank(\sum_t P_{V_t})$ within $5\%$.
  Passed Day~8 for $T=2$, $r\in\{4,8,16,32\}$.
\end{itemize}

\paragraph{Day-10 verification (general $H_t$, common $D_0$).}
Numerics in \texttt{code/synthetic/day10\_general\_ht.py} verify
Theorem~\ref{thm:general} and Lemma~\ref{lem:floor-general} under four
eigenvalue specs for a common $D_0$ (isotropic, geometric, linear,
two-scale) across $T\in\{2,3,4\}$ and $r\in\{4,8\}$: floor and
$\E[\|\eta\|^2]$ match theory to five digits, MP-aware rotational-
invariance test passes at $n \in \{3\times 10^4, 10^5\}$, and the
$\log_2\text{excess}_{\bar H}$ slope is $-2.00\pm 0.01$ under $\eta$-
based adaptive quantization.

\paragraph{Day-11 verification (task-dependent $D_t$).} Numerics in
\texttt{code/synthetic/day11\_task\_dep\_Dt.py} verify the same
Theorem~\ref{thm:general} and Lemma~\ref{lem:floor-general} under
\emph{task-mismatched} per-task spectra (eight $(D_t)_t$
configurations including \texttt{iso/geom}, \texttt{geom/twoscale},
\texttt{iso/geom/twoscale/lin} across $T\in\{2,3,4\}$). Floor and
$\E[\|\eta\|^2]$ match theory to four--five digits; the MP-aware
rotational-invariance spread is below the isotropic UB in all four
cells at $n\in\{3\times 10^4, 10^5\}$; $\log_2\text{excess}_{\bar H}$
slope is $-2.00\pm 0.01$ across all six shared/random cells. The
empirical evidence confirms the $D_t$-free character of the LB
argued in Step 5 and the remark on ``why the bound doesn't depend
on $(D_t)_t$.''

\section{Matching achievability (Stiefel floor-zero regime)}\label{sec:achv}

Phase 2 target: an explicit encoder--decoder pair matching the lower
bound of Theorem~\ref{thm:general} up to a universal multiplicative
constant. We give the construction for the floor-zero regime
($V_t$ Stiefel-random iid, $Tr\leq d$, so $\deff=Tr$ a.s.\ under
$P^\star$), which is the generic LoRA case. The general regime
(floor $>0$, shared/overlapping $V_t$) is addressed in
Remark~\ref{rem:sharedv-gap} below.

\paragraph{Algorithm (Gaussian-QR + uniform scalar quantization).}
Encoder and decoder share randomness (a seed indexing an orthogonal
rotation $O\in O(\deff)$ drawn uniformly from Haar measure).
\begin{enumerate}
  \item Encoder computes $\Hbar$, $\tauH$, an orthonormal basis $Q$
  of $\range(\Hbar)$ with eigenvalues $w_1,\dots,w_\deff$, and
  $\eta := \Hbar^{1/2}\tauH$ represented in the $Q$-basis as
  $\eta_Q = (\sqrt{w_i}\,\langle Q_i,\tauH\rangle)_{i=1}^\deff$.
  \item Apply the rotation: $\tilde\eta := O\eta_Q \in \R^\deff$.
  \item Uniform scalar-quantize each coord of $\tilde\eta$ in
  $[-c\sigma_{\mathrm{pc}},\, c\sigma_{\mathrm{pc}}]$ with $b$ bits,
  where $\sigma_{\mathrm{pc}} = B/\sqrt{Tr}$ and $c$ is a clipping
  constant (e.g.\ $c = 5$). Total rate $R = b\deff$.
  \item Encoder transmits the quantization indices. Decoder
  unquantizes to $\tilde\eta_q$, applies $O^\top$ to get $\eta_{Q,q}$,
  and emits $w^\star := Q\,\mathrm{diag}(1/\sqrt{w_i})\,\eta_{Q,q}
  = \Hbar^{+1/2}\eta_q$ (on $\range(\Hbar)$).
\end{enumerate}

\begin{theorem}[Matching achievability, floor-zero Stiefel]\label{thm:achv}
Under $P^\star$ with Stiefel-random $V_t$ and $Tr\leq d$, for any
fixed positive-definite $(D_t)_{t=1}^T$, the above algorithm at
rate $R$ bits attains
\begin{equation}\label{eq:thm-achv}
  \E_{P^\star}\!\left[\max_t(w^\star-\tau_t)^\top H_t(w^\star-\tau_t)\right]
  \;\leq\; C\cdot B^2\cdot 2^{-2R/(Tr)}
\end{equation}
for a universal constant $C = T c^2 / 3$. Combined with the
specialization of Theorem~\ref{thm:general} to this regime,
\[
  \E_{P^\star}\!\left[\max_t(w^\star-\tau_t)^\top H_t(w^\star-\tau_t)\right]
  \;=\; \Theta\!\left(B^2\cdot 2^{-2R/(Tr)}\right),
\]
i.e.\ the Bayesian rate--distortion function of merging under
$P^\star$ is pinned down up to constants in the Stiefel floor-zero
regime.
\end{theorem}

\begin{proof}
Five steps.

\emph{Step 1 (zero-floor identity).} Under $P^\star$ with $Tr\leq d$,
Lemma~\ref{lem:floor-general} gives
$\E_{P^\star}[T^{-1}\sum_t(\tau_t-\tauH)^\top H_t(\tau_t-\tauH)] = 0$.
Since each summand is nonnegative, $(\tau_t-\tauH)^\top H_t
(\tau_t-\tauH) = 0$ a.s.\ for every $t$, equivalently $H_t\tauH =
H_t\tau_t$ a.s. (``perfect per-task reconstruction at $\tauH$'').

\emph{Step 2 (distortion reduces to quantization error on $\eta$).}
Let $\delta w := w^\star-\tauH\in\range(\Hbar)$. Then
\begin{align*}
  (w^\star-\tau_t)^\top H_t(w^\star-\tau_t)
  &= ((\tauH-\tau_t)+\delta w)^\top H_t((\tauH-\tau_t)+\delta w) \\
  &= 0 + 2(\delta w)^\top H_t(\tauH-\tau_t) + (\delta w)^\top H_t(\delta w) \\
  &= (\delta w)^\top H_t(\delta w),
\end{align*}
using Step 1 ($H_t(\tauH-\tau_t) = 0$) for both the zero and the
cross term.

\emph{Step 3 (max $\leq T\cdot$ avg).} Summing over $t$:
$\sum_t (\delta w)^\top H_t \delta w = T(\delta w)^\top\bar H \delta w$,
so
\[
  \max_t (\delta w)^\top H_t \delta w \;\leq\;
  \sum_t (\delta w)^\top H_t \delta w \;=\;
  T\,(\delta w)^\top\bar H \delta w \;=\; T\,\|\bar H^{1/2}\delta w\|^2
  \;=\; T\,\|\delta\eta\|^2,
\]
where $\delta\eta := \eta_q - \eta$ is the quantization error in
$\eta$-space.

\emph{Step 4 (quantization error under Gaussian-QR mixing).} After
the orthogonal rotation $O$, $\tilde\eta = O\eta_Q$ has per-coord
second moment $\E[\tilde\eta_i^2] = \E\|\eta\|^2/\deff = B^2/(Tr)$
by Lemma~\ref{lem:floor-general}. Uniform scalar quantization with
range $[-c\sigma_{\mathrm{pc}},c\sigma_{\mathrm{pc}}]$ and $b$ bits
has per-coord MSE at most $\Delta^2/12 = c^2\sigma_{\mathrm{pc}}^2
2^{-2b}/3$, giving
\[
  \E\|\delta\tilde\eta\|^2 \;\leq\; \deff\cdot\frac{c^2 B^2}{3Tr}\cdot 2^{-2b}
  \;=\; \frac{c^2 B^2\deff}{3Tr}\cdot 2^{-2R/\deff}.
\]
Since $O$ is orthogonal, $\|\delta\eta\|^2 = \|\delta\tilde\eta\|^2$;
the same bound holds in $\eta$-space. In the Stiefel regime
$\deff = Tr$, so
$\E\|\delta\eta\|^2 \leq (c^2 B^2/3)\cdot 2^{-2R/(Tr)}$.

\emph{Step 5 (combine).} From Step 3 and Step 4,
\[
  \E_{P^\star}[\max_t(w^\star-\tau_t)^\top H_t(w^\star-\tau_t)]
  \;\leq\; T\,\E\|\delta\eta\|^2
  \;\leq\; \frac{T c^2}{3}\cdot B^2\cdot 2^{-2R/(Tr)}.
\]
This is \eqref{eq:thm-achv} with $C = Tc^2/3$.
\end{proof}

\begin{remark}[Empirical constant]
\texttt{code/synthetic/day12\_achievability.py} and
\texttt{day13\_achievability\_fixes.py} verify Theorem~\ref{thm:achv}
numerically: the ratio $\mathrm{UB}_{\text{empirical}} /
\mathrm{LB}_{\text{Thm 8, }c_{\mathrm{TQ}}=1}$ is
$\leq 13$ across $T\in\{2,3\}$, $r=4$, $d=128$, iso+iso, iso+geom,
and iso+geom+twoscale task spectra --- constant across all five rate
points tested ($b\in\{2,4,6,8,10\}$). The empirical constant is
consistent with $C = Tc^2/3 \approx 17$ for $c=5, T=2$, and with
TurboQuant's $c_{\mathrm{TQ}}\approx 1$.
\end{remark}

\begin{remark}[Shared-$V$ / floor $>0$ gap]\label{rem:sharedv-gap}
Theorem~\ref{thm:achv} is stated for the floor-zero regime. In
shared-$V$ (or more generally non-zero-floor) regimes, quantization
of $w^\star$ around the Chebyshev center $w_{\text{cheb}}$ via
$H_t$-weighted KKT analysis shows: at $w_{\text{cheb}}$ with $|A|$
active tasks, the active gradients $\{g_t = H_t(w_{\text{cheb}}-\tau_t)
: t\in A\}$ span $|A|-1$ dimensions in $\range(\Hbar)$. Any explicit
linear encoder that quantizes $w_{\text{cheb}}$ with zero-mean noise
$\delta w$ produces
$\max_t\,(w^\star-\tau_t)^\top H_t(w^\star-\tau_t) \approx
\text{cheb}^2 + 2\max_{t\in A}|g_t^\top \delta w| + \|\delta w\|^2$.
The linear term is $O(\|g\|\cdot\sigma_\parallel)$ with
$\sigma_\parallel$ the quantization stdev in the parallel direction,
and for uniform bit allocation this decays only as $2^{-R/\deff}$,
giving slope $-1$ (not $-2$) on $\log_2$ excess vs $b = R/\deff$.
A balanced allocation $R_\parallel = 2R(|A|-1)/(\deff+|A|-1)$,
$R_\perp = R\deff/(\deff+|A|-1) - \deff$ (after dimensional
accounting) improves the excess to $\Theta(2^{-2R/(\deff+|A|-1)})$,
which is strictly better than naive $2^{-R/\deff}$ but strictly
worse than Theorem~\ref{thm:general}'s LB exponent $2R/\deff$ for
$|A|>1$.
Empirical verification at $T=|A|=2$, $r=\deff=4$: slope improves
from $-1.12$ (naive Day~13 algorithm,
\texttt{code/synthetic/day13\_achievability\_fixes.py}) to $-1.66$
(null-space split, iso+iso,
\texttt{code/synthetic/day14\_shared\_v\_gap.py}), matching the
theoretical prediction $-2\deff/(\deff+|A|-1) = -8/5 = -1.60$ to
within sampling noise. The Day-14 algorithm initially failed on
non-iso $D_t$ cells (slope $\approx 0$) due to a smooth-max
Newton-solver convergence bug; we fixed this with a closed-form
$T=2$ Chebyshev solver that roots-finds the KKT equation
$D_1(w(\alpha)) = D_2(w(\alpha))$ for
$w(\alpha) = (\alpha H_1 + (1-\alpha) H_2)^{-1}(\alpha H_1 \tau_1
+ (1-\alpha) H_2 \tau_2)$ via Brent's method
(\texttt{day14b\_cheb\_T2\_closedform.py}); the residual $|D_1-D_2|$
lands at $10^{-9}$ (machine precision for a 1-D root-find). Combined
with fractional per-coord bit widths on the perpendicular subspace
(\texttt{day14c\_fractional\_bits.py}) and a locked scalar clip
$c = 11.5\sigma_{\mathrm{pc}}$
(\texttt{day14g\_final\_lock.py}, $n_{\mathrm{trials}}=1000$), the
null-space split achieves slope $-1.60 \pm 0.10$ across five
anisotropy regimes: iso+iso $-1.704$, iso+geom $-1.604$, geom+twos
$-1.502$, lin+twos $-1.591$, geom+iso $-1.603$ (bootstrap 95\% CI
$\pm 0.010$ on each). The $\pm 0.10$ spread reflects anisotropy
structure (iso cells track closer to the scalar-quant limit $-2$;
geom+twos undershoots due to $H_t$-metric distortion on the perp
basis), not sampling noise. Closing the remaining gap to LB's
$-2$ likely requires non-linear (random coding) encoders; this is
an open question at the intersection of rate--distortion theory and
combinatorial optimization. Flagged as Phase~2.5 item in
\S\ref{sec:open}.

\paragraph{Phase 2.5 Day 15--16 update.}
Two follow-up experiments pinned down the structure of the gap.
\emph{Day 15}: the null-space-split algorithm was extended to $T\geq 3$
using a cvxpy SOCP Chebyshev-center solver with a post-solve
Gauss--Newton KKT refinement
(\texttt{code/synthetic/day15\_cheb\_general\_T.py}). KKT residuals
land at $10^{-15}$ on the active set; for $T=2,3,4$ and $r=4$, the
empirical slope of $\max_t D_t(w^\star) - \mathrm{cheb}^2$ (excess
over the \emph{per-trial} Chebyshev value, not the average floor,
which is the quantity the theory bounds) matches or beats the linear
prediction $-2r/(r + |A| - 1)$ across all anisotropy regimes: e.g.\
$T=3$ iso+geom+twos slope $-1.62$ vs predicted $-1.40$, $T=4$ iso+4
slope $-1.42$ vs predicted $-1.30$. (The subtraction of cheb$^2$ is
essential: excess-over-avg-floor saturates at
$\mathrm{cheb}^2 - \E[\mathrm{floor}] \neq 0$ for $T\geq 3$ under
anisotropy, which is why Day 14's early attempt appeared to fail.)
\emph{Day 16}: a valid rate-$R$ random-codebook encoder
(\texttt{day16\_sharpened\_lb\_check.py}, $2^R$ iid Gaussian codewords
at origin, decoder does nearest-neighbour in the $\max_t D_t$ metric)
was run at $T=2, r=3$ and $T=3, r=3$ shared-$V$. The random-codebook
slopes \emph{strictly beat the linear null-space prediction} at
$T=3$: $-1.45$ vs linear $-1.20$. This \emph{rules out} sharpening
the LB to match the linear UB: if the LB were $-2r/(r+|A|-1) = -1.20$,
the random-codebook encoder's slope $-1.45$ would violate it. Hence
the true exponent in the shared-$V$/floor-positive regime lies
\emph{strictly between} $-2r/r$ (the current Theorem~\ref{thm:general} LB,
loose here) and $-2r/(r+|A|-1)$ (the linear-encoder UB, also loose).
Closing this gap is a genuine open question at the intersection of
rate--distortion theory and combinatorial optimization and is beyond
Phase 2's explicit-construction scope.
\end{remark}

\section{Remaining open items for Phase 1}\label{sec:open}

Items 1--3 of the Day-8 list (prove the rank-$r$ RD lower bound,
closed form for the floor, explicit $\deff$) are resolved by
Theorem~\ref{thm:rankr} and Lemma~\ref{lem:floor-rankr} above. Day-8
item 4 (functional distortion / general $H_t$) is resolved in
\S\ref{sec:general-ht} under the \emph{full} task-dependent-spectrum
regime $(D_t)_{t=1}^T$ (Theorem~\ref{thm:general}). On Day 11 the
Day-10 hedge about ``deterministic mismatched $D_t$'s break Step~5''
was retired: close re-reading + numerical verification
(\texttt{code/synthetic/day11\_task\_dep\_Dt.py}) showed that the
argument is $D_t$-free at both the floor step (per-task identity
$H_t H_t^+ H_t = H_t$) and the equivariance step ($v_t$ is
$O(d)$-free regardless of which $D_t$ it uses). Phase 1 is therefore
fully closed. Remaining items:
\begin{enumerate}
  \item \textbf{Robustness of $\tauH$ for small-overlap $V_t$.} When
  $V_t$'s are nearly orthogonal, $\Hbar$ is close to $(r/d)\cdot
  I_{\sum V_t}$ and $\tauH$ is stable; when they overlap irregularly,
  $\Hbar$ can be ill-conditioned. The current
  Theorems~\ref{thm:rankr}--\ref{thm:general} bounds degrade
  gracefully via $\deff$, but the $O(\deff)$ stabilizer argument in
  Step~5 relies on $\deff$ being a.s.\ constant under $P$/$P^\star$,
  which holds for Stiefel-random $V_t$ but may fail for more general
  overlap distributions. A quantitative conditioning assumption may
  be needed for the general case.
  \item \textbf{Achievability (Phase 2).} Closed in the floor-zero
  Stiefel regime by Theorem~\ref{thm:achv}. The shared-$V$ /
  floor $>0$ regime has a genuine exponent gap
  (Remark~\ref{rem:sharedv-gap}): linear null-space-split encoders
  achieve $2^{-2R/(r+|A|-1)}$, random-codebook (non-linear) encoders
  strictly beat this empirically, and the current LB
  $2^{-2R/r}$ is loose here. Phase~2.5 investigation (Days 15--16)
  extended the linear construction to $T\geq 3$ via SOCP and
  empirically ruled out sharpening the LB to match the linear UB.
  The exact RD exponent for the shared-$V$ regime is open and
  beyond Phase 2's scope.
  \item \textbf{Real-LLM empirical validation.} Do the Theorem~8
  LB and Theorem~9 UB hold on real LoRA fine-tunes of LLMs, not
  just Stiefel-random synthetic tuples? Phase~3 item.
\end{enumerate}

\end{document}
